\subsection{Labor heterogeneity and reallocation}
\label{app:rewrite-labor-heterogeneity}

\paragraph{Joint demand and global monopoly supply.}
Write $H=AN$ and $H_j=AL_j$. At fixed positive sectoral labor supplies,
define the optimized composite at total AI supply $X$ by
\[
 G(X)=\max_{x_S+x_C=X,\ x_j\geq0}
 \{Z_S(H_S,x_S)+Z_C(H_C,x_C)\}.
\]
Capital allocation in equation~\eqref{eq:rewrite-heterogeneous-capital}
then gives $Y=K^\alpha G(X)^{1-\alpha}$.
Both CES functions are concave, and their marginal products of AI decrease
strictly on the positive domain. The allocation is therefore unique.
Let $\mu=G'(X)$ denote their common marginal product when both use AI.
Set
\[
 c_S=\omega_{XS}^{\sigma_S/(\sigma_S-1)}<1,\qquad
 b_C=\omega_{XC}^{\sigma_C/(\sigma_C-1)}>1.
\]
Solving $\partial Z_j/\partial x_j=\mu$ yields
\begin{equation}
 x_j(\mu)=H_j\left[
 \frac{\omega_{Lj}}{(\mu/\omega_{Xj})^{\sigma_j-1}-\omega_{Xj}}
 \right]^{\sigma_j/(\sigma_j-1)}.
 \label{eq:rewrite-heterogeneous-allocation}
\end{equation}
For $S$, use this expression for every $\mu>c_S$. For $C$, use it only
when $\mu<b_C$, and set $x_C=0$ otherwise. Total supply decreases
continuously from infinity to zero as $\mu$ runs from $c_S$ to infinity.
The resulting inverse demand is
\[
 P(X)=(1-\alpha)[K/G(X)]^\alpha G'(X).
\]
It is strictly decreasing. Its elasticity must account for the changing
allocation of services, rather than hold sectoral capital fixed.

Operating profit $P(X)X-X/B$ has a positive finite global maximum.
At small $X$, only $S$ uses AI, and revenue is proportional to
$X^{1-1/\sigma_S}$ to leading order. This dominates compute cost $X/B$.
At large $X$, equation~\eqref{eq:rewrite-heterogeneous-allocation} gives
$x_C/H$ a finite limit, $x_S/X\to1$, $G(X)/X\to c_S$, and
$G'(X)\to c_S$. Revenue is then of order $X^{1-\alpha}$, whereas cost
is linear. Profit is continuous, vanishes at zero and tends to minus infinity
at infinity, establishing the claim. This argument does not assert uniqueness
at all states, so the developer must select a global maximum.

For the limiting supply, let $k=K/H$ and $G(X)=H\mathcal G(X/H)$.
At a fixed supply-to-capital ratio $y=X/K>0$, profit divided by $K$ converges,
as $k\to\infty$ and $B\to B_\infty$, to
\begin{equation}
 f_\infty(y)=(1-\alpha)c_S^{1-\alpha}y^{1-\alpha}-y/B_\infty.
 \label{eq:rewrite-heterogeneous-profit-limit}
\end{equation}
To verify uniform convergence on a compact positive interval of $y$,
expand the substituting CES at large AI use. With
$\nu=(\sigma_S-1)/\sigma_S>0$, the expansion and the bounded complementary
allocation give $\mathcal G(x)=c_Sx+O(x^{1-\nu})$ and
$\mathcal G'(x)=c_S+O(x^{-\nu})$. For $y$ bounded away from zero,
these errors vanish uniformly after setting $x=ky$ and normalizing profit.
All global
maximizers lie in such an interval for large $k$. To verify this last step,
use concavity to obtain $x\mathcal G'(x)\leq\mathcal G(x)$ and a supporting
bound $\mathcal G(x)\leq a+bx$ with finite positive constants. Normalized
profit is bounded above by
\[
 (1-\alpha)(a/k+by)^{1-\alpha}-y/B.
\]
This excludes arbitrarily large $y$. It also shows that $y\to0$ gives
limiting normalized profit at most zero, whereas a fixed positive $y$ can
give positive profit in equation~\eqref{eq:rewrite-heterogeneous-profit-limit}.
The limit objective has a unique maximum. Uniform convergence on the
remaining compact interval therefore implies, for every global selection,
\begin{equation}
 \frac{X}{K}\to
 [B_\infty(1-\alpha)^2c_S^{1-\alpha}]^{1/\alpha},\qquad
 p_X\to p_*\equiv\frac1{B_\infty(1-\alpha)}.
 \label{eq:rewrite-heterogeneous-supply-limit}
\end{equation}
This proof requires no assumption of global concavity of the finite-state
monopoly objective.

\begin{proof}[Proof of Proposition~\ref{prop:rewrite-heterogeneous-wages}]
The preceding supply limit and $G/X\to c_S$ give
$r\to\mathcal R_S(B_\infty)$ and $Y/K$ a positive finite limit.
The household Euler equation implies average consumption-per-person growth
$\mathcal R_S(B_\infty)-\rho$. Since $C/K$ converges to a positive
constant, capital and output per worker have the same average growth $g$.

Equation~\eqref{eq:rewrite-heterogeneous-allocation} gives $X_C/H_C$
a finite positive limit because $c_S<b_C$. Hence $Z_C/H_C$ converges to
a positive constant. Also $Q\to Q_*=p_*/c_S>0$.
The wage equation~\eqref{eq:rewrite-heterogeneous-prices} then gives
$w_C/A$ a positive finite limit, and therefore $g_{w_C}=\gamma$.
Since $Z_C=O(H)$ and $G/H\to\infty$, $Z_C/G\to0$ and
$Z_S/G\to1$. Thus $Z_S$ has average growth $n+g$.
Taking logarithms of the wage equation for $S$ gives
\[
 g_{w_S}=\gamma+\frac{g-\gamma}{\sigma_S}.
\]
Both wages grow more slowly than output per worker, and their employment
shares are fixed. Each labor-income share in aggregate output therefore
vanishes. Since capital is allocated proportionally to composites,
$Y_C/Y=Z_C/G\to0$. The positive limiting ratios in activity $C$ imply
that it remains active at sufficiently late dates, despite its vanishing share.
For competitive provision, replace the monopoly condition with $P(X)=1/B_0$.
The same allocation argument gives the stated marginal-cost return and wages.
\end{proof}

\paragraph{A finite supply threshold for mobile labor.}
Let
\[
 b_L=\omega_{LC}^{\sigma_C/(\sigma_C-1)}.
\]
For every feasible pair $(H_C,x_C)$, complementarity gives
$Z_C(H_C,x_C)\leq b_L H_C$.
The marginal product of effective labor in $S$, evaluated using all
effective labor and all AI there, is
\[
 Z_{S,H}(H,X)=\omega_{LS}
 [\omega_{LS}+\omega_{XS}(X/H)^{(\sigma_S-1)/\sigma_S}]^{1/(\sigma_S-1)}.
\]
It exceeds $b_L$ at every $X/H$ above the finite threshold
\begin{equation}
 x_{\mathrm{mob}}=
 \left[
 \frac{(b_L/\omega_{LS})^{\sigma_S-1}-\omega_{LS}}{\omega_{XS}}
 \right]^{\sigma_S/(\sigma_S-1)}.
 \label{eq:rewrite-mobile-threshold}
\end{equation}
For any alternative allocation with $H_C>0$, monotonicity in AI and
concavity in labor imply
\begin{align*}
 Z_S(H-H_C,X-x_C)+Z_C(H_C,x_C)
 &\leq Z_S(H-H_C,X)+b_LH_C\\
 &\leq Z_S(H,X)-[Z_{S,H}(H,X)-b_L]H_C.
\end{align*}
Above the threshold, the final expression is strictly below $Z_S(H,X)$.
An allocation with $H_C=0$ and $x_C>0$ also wastes services, because the
complementary activity produces zero without labor. Thus the unique
efficient allocation above this sufficient threshold assigns labor and AI
entirely to $S$, and optimal capital allocation assigns all capital there too.

\begin{proof}[Proof of Corollary~\ref{cor:rewrite-mobile-labor}]
With mobile labor, maximize $Z_S+Z_C$ over both labor and AI allocation.
This frontier is concave and homogeneous in total effective labor and AI.
By the finite threshold above, it equals $Z_S(H,X)$ whenever
$X/H>x_{\mathrm{mob}}$. Its high-supply slope is therefore $c_S$,
and it admits a supporting affine bound per unit of $H$.
The compactness and limiting-profit argument leading to
equation~\eqref{eq:rewrite-heterogeneous-supply-limit} applies to this
frontier as well. Hence any global monopoly choice has $X/K$ converging
to the same strictly positive constant. Since $K/H\to\infty$, $X/H$
eventually exceeds $x_{\mathrm{mob}}$ and stays above it.
There is therefore a finite date after which $L_C=X_C=K_C=0$.
Production and the dated input prices then coincide with the original
single-activity technology using $\sigma_S,\omega_{LS},\omega_{XS}$.
The limiting return and the household Euler equation give $g$ as stated,
and the wage equation with $L_S=N$ gives the claimed average wage growth.
This identifies dated allocations and limiting rates along the assumed
equilibrium. It does not transfer an existence theorem or a global
research-verification argument from the original model.
\end{proof}
